% GENERATED FILE. Edit canonical lessons, then run npm run generate:books.
% canonical-source-hash: fnv1a64:2867f06784be7547
% canonical-lessons: AR-C117-tilfaz, AR-C117-jarida, AR-C117-majalla, AR-C117-risala, AR-C117-zarf

\chapter{Television, Newspaper, Magazine, Letter, Envelope}
\label{ch:ar-television-newspaper-magazine-letter-envelope}
\hlchaptermodality{117}

\begin{chapteropening}
\textbf{By the end of this chapter:} \emph{I can say a television, a newspaper, a magazine, a letter and an envelope.}

Complete the last lesson of chapter 117: I can say a television, a newspaper, a magazine, a letter and an envelope.
\end{chapteropening}

\section[\texorpdfstring{\ar{تلفاز} (tilfāz)}{تلفاز (tilfāz)}]{\ar{تلفاز} (tilfāz) — a television}

\label{lesson:AR-C117-tilfaz}



\begin{quote}
Before the new one: say the Arabic for a fan, then the Arabic for a computer.
\end{quote}

\subsection*{You'll want to know: \ar{تلفاز}}
\textbf{\ar{تلفاز}} — \emph{tilfāz} — \textquotedblleft{}a television\textquotedblright{}.

\begin{cousinweb}
\textbf{\ar{تلفاز}} (\emph{tilfāz}) is a modern word borrowed from European \emph{television} and fitted to an Arabic pattern; it has no Arabic root.
\end{cousinweb}

\begin{grammarlens}[title={What you've built}]
One more word for this chapter.
\end{grammarlens}

\subsection*{Your turn}
\begin{itemize}
\raggedright
  \item \emph{Say it:} \emph{tilfāz}
  \item \emph{Say it:} \emph{tilfāz}, once more
  \item \emph{Say it:} say \emph{tilfāz}
  \item \emph{Recall:} say the Arabic for a fan, then the Arabic for a computer, then say \emph{tilfāz} again
\end{itemize}

\begin{tcolorbox}[breakable,colback=teal!4,colframe=teal!35!black,title={Before you move on}]
What does \ar{تلفاز} mean? (\textquotedblleft{}A television\textquotedblright{}.) Say it once more.
\end{tcolorbox}
\section[\texorpdfstring{\ar{جريدة} (jarīda)}{جريدة (jarīda)}]{\ar{جريدة} (jarīda) — a newspaper}

\label{lesson:AR-C117-jarida}



\begin{quote}
Before the new one: say the Arabic for a computer, then the Arabic for a television.
\end{quote}

\subsection*{You'll want to know: \ar{جريدة}}
\textbf{\ar{جريدة}} — \emph{jarīda} — \textquotedblleft{}a newspaper\textquotedblright{}.

\begin{cousinweb}
\textbf{\ar{جريدة}} (\emph{jarīda}) is built on the native Arabic root \emph{j-r-d}, \textquotedblleft{}to strip\textquotedblright{} (Lane). Lane's \emph{jarīda} is a palm stalk stripped of its leaves, and also a register; the newspaper is a modern sense.
\end{cousinweb}

\begin{grammarlens}[title={What you've built}]
One more word for this chapter.
\end{grammarlens}

\subsection*{Your turn}
\begin{itemize}
\raggedright
  \item \emph{Say it:} \emph{jarīda}
  \item \emph{Say it:} \emph{jarīda}, once more
  \item \emph{Say it:} say \emph{jarīda}
  \item \emph{Recall:} say the Arabic for a computer, then the Arabic for a television, then say \emph{jarīda} again
\end{itemize}

\begin{tcolorbox}[breakable,colback=teal!4,colframe=teal!35!black,title={Before you move on}]
What does \ar{جريدة} mean? (\textquotedblleft{}A newspaper\textquotedblright{}.) Say it once more.
\end{tcolorbox}
\section[\texorpdfstring{\ar{مجلة} (majalla)}{مجلة (majalla)}]{\ar{مجلة} (majalla) — a magazine}

\label{lesson:AR-C117-majalla}



\begin{quote}
Before the new one: say the Arabic for a television, then the Arabic for a newspaper.

Say \textquotedblleft{}the weather,\textquotedblright{} then \textquotedblleft{}It is hot.\textquotedblright{} (\emph{Aṭ-ṭaqs}; \emph{al-jaw ḥārr} — no verb \textquotedblleft{}to be.\textquotedblright{})
\end{quote}

\subsection*{You'll want to know: \ar{مجلة}}
\textbf{\ar{مجلة}} — \emph{majalla} — \textquotedblleft{}a magazine\textquotedblright{}.

\begin{cousinweb}
\textbf{\ar{مجلة}} (\emph{majalla}) is filed by Lane under \emph{j-l-l} as a book or scroll, and one of his sources ties it to \emph{jill}, \textquotedblleft{}a book's cover\textquotedblright{}. It is often compared with Aramaic \emph{megillā}, \textquotedblleft{}scroll\textquotedblright{}, and may be borrowed from it; the magazine is a modern sense.
\end{cousinweb}

\begin{grammarlens}[title={What you've built}]
One more word for this chapter.
\end{grammarlens}

\subsection*{Your turn}
\begin{itemize}
\raggedright
  \item \emph{Say it:} \emph{majalla}
  \item \emph{Say it:} \emph{majalla}, once more
  \item \emph{Say it:} say \emph{majalla}
  \item \emph{Recall:} say the Arabic for a television, then the Arabic for a newspaper, then say \emph{majalla} again
\end{itemize}

\begin{tcolorbox}[breakable,colback=teal!4,colframe=teal!35!black,title={Before you move on}]
What does \ar{مجلة} mean? (\textquotedblleft{}A magazine\textquotedblright{}.) Say it once more.
\end{tcolorbox}
\section[\texorpdfstring{\ar{رسالة} (risāla)}{رسالة (risāla)}]{\ar{رسالة} (risāla) — a letter, a message}

\label{lesson:AR-C117-risala}



\begin{quote}
Before the new one: say the Arabic for a newspaper, then the Arabic for a magazine.
\end{quote}

\subsection*{You'll want to know: \ar{رسالة}}
\textbf{\ar{رسالة}} — \emph{risāla} — \textquotedblleft{}a letter, a message\textquotedblright{}.

\begin{cousinweb}
\textbf{\ar{رسالة}} (\emph{risāla}) is built on the native Arabic root \emph{r-s-l}, \textquotedblleft{}to send\textquotedblright{} (Lane). \emph{Risāla} is a thing sent: a message, a letter.
\end{cousinweb}

\begin{grammarlens}[title={What you've built}]
One more word for this chapter.
\end{grammarlens}

\subsection*{Your turn}
\begin{itemize}
\raggedright
  \item \emph{Say it:} \emph{risāla}
  \item \emph{Say it:} \emph{risāla}, once more
  \item \emph{Say it:} say \emph{risāla}
  \item \emph{Recall:} say the Arabic for a newspaper, then the Arabic for a magazine, then say \emph{risāla} again
\end{itemize}

\begin{tcolorbox}[breakable,colback=teal!4,colframe=teal!35!black,title={Before you move on}]
What does \ar{رسالة} mean? (\textquotedblleft{}A letter, a message\textquotedblright{}.) Say it once more.
\end{tcolorbox}
\section[\texorpdfstring{\ar{ظرف} (ẓarf)}{ظرف (ẓarf)}]{\ar{ظرف} (ẓarf) — an envelope}

\label{lesson:AR-C117-zarf}



\begin{quote}
Before the new one: say the Arabic for a magazine, then the Arabic for a letter.
\end{quote}

\subsection*{You'll want to know: \ar{ظرف}}
\textbf{\ar{ظرف}} — \emph{ẓarf} — \textquotedblleft{}an envelope\textquotedblright{}.

\begin{cousinweb}
\textbf{\ar{ظرف}} (\emph{ẓarf}), \textquotedblleft{}a vessel, a receptacle\textquotedblright{}, is an old native noun, filed under the root \emph{ẓ-r-f} (Lane), whose verb means \textquotedblleft{}to be elegant, clever\textquotedblright{}. Lane's \emph{ẓarf} is a vessel; the envelope is a later sense.
\end{cousinweb}

\begin{grammarlens}[title={What you've built}]
One more word for this chapter.
\end{grammarlens}

\subsection*{Your turn}
\begin{itemize}
\raggedright
  \item \emph{Say it:} \emph{ẓarf}
  \item \emph{Say it:} \emph{ẓarf}, once more
  \item \emph{Say it:} say \emph{ẓarf}
  \item \emph{Recall:} say the Arabic for a magazine, then the Arabic for a letter, then say \emph{ẓarf} again
\end{itemize}

\begin{tcolorbox}[breakable,colback=teal!4,colframe=teal!35!black,title={Before you move on}]
What does \ar{ظرف} mean? (\textquotedblleft{}An envelope\textquotedblright{}.) Say it once more.
\end{tcolorbox}
